La primera salida de la IA respetó la sintaxis \texttt{sequenceDiagram} y declaró correctamente el actor \texttt{U as Administradora}, pero presentó las siguientes deficiencias respecto a la arquitectura del proyecto:
\begin{itemize}[leftmargin=*,nosep]
\item \textbf{Omisión del flujo de autenticación contra base de datos:} La IA asumió que la administradora ya disponía de un token JWT activo y omitió el login. Para mantener la coherencia con el caso de uso individual (\texttt{cu\_us05}, que incluye \texttt{Autenticar administradora}) y con el diagrama de clases, se incorporó \texttt{POST /api/auth/login (email, password)}, la consulta \texttt{findUsuarioByEmail(email)} y el método \texttt{Usuario.iniciarSesion(email, password)}.
\item \textbf{Rutas REST y esquema de base de datos en inglés:} La IA empleó \texttt{GET /api/admin/orders?paymentStatus=PAID} y \texttt{PATCH /api/admin/orders/\{id\}/status}. Se corrigieron por las rutas del glosario común: \texttt{GET /api/admin/pedidos} y \texttt{PATCH /api/admin/pedidos/idPedido (estado ENVIADO)}.
\item \textbf{Uso de los métodos de las clases \texttt{Administradora} y \texttt{Pedido}:} En lugar de sentencias en prosa o SQL arbitrario, se especificaron las llamadas a los métodos definidos en \texttt{clases.mmd}: \texttt{Administradora.listarPedidosPagados()}, \texttt{Administradora.marcarPedidoEnviado(idPedido)} y \texttt{Pedido.cambiarEstado(ENVIADO)}, eliminando la rama \texttt{404} redundante y añadiendo \texttt{mirrorActors: false}.
\end{itemize}
